In this section, we prove our main theorems, \Cref{thm:main-A} and \Cref{thm:main-B}. We do this by first proving the most immediate directions of the equivalences in \Cref{thm:main-A}. We then prove \Cref{thm:main-B} in several lemmas and then deduce the remaining directions of \Cref{thm:main-A} as a corollary.

\subsection{Proofs of some parts of \texorpdfstring{\Cref{thm:main-A}}{Theorem A}}

In this subsection, we prove that implications in \Cref{thm:main-A} that follow easily from work of Gasharov--Reiner~\cite{GasharovReiner}, namely \Cref{cond:A2} $\implies$ \Cref{cond:A3} $\implies$ \Cref{cond:A1}. We complete the final implication \Cref{cond:A1} $\implies$ \Cref{cond:A2} by directly proving the more general statement that \Cref{cond:B1} $\implies$ $\calW\subseteq S_n(\calP)$ in \Cref{prop:1-implies-avoidance}.


\begin{lemma}[\Cref{cond:A3} implies \Cref{cond:A1}]
\label{lem: A3 implies A1}
    If $X_w$ is defined by the inclusions stated in \Cref{cond:A3}, then either $w=\id$ and $J_\id=(x_1,\dots, x_n)$ or $w\neq \id$ and $J_w$ is generated by the following squarefree monomials, where $i = i(w)$ and $j=j(w)$,
    \begin{itemize}
        \item The variables $x_1,\dots,x_{i-2}$, and $x_{i-1}$,
        \item The products $x_i\cdots x_{w^{-1}(i)}$ and $x_{w^{-1}(j)}\cdots x_j$, and
        \item The variables $x_{j+1},\dots,x_{n-1}$, and $x_n$.
    \end{itemize}
\end{lemma}

\begin{proof}
    If $w=\id$, then as ideals in $R_n$,
    \[
    J_w = (\fraks_{s_i}:1\leq i<n) = (x_1+x_2+\cdots + x_i:1\leq i<n)= (x_1,\dots, x_n).
    \]
    Otherwise, suppose $w\neq \id$ and $X_w$ is defined by the inclusion stated in \Cref{cond:A3}. Applying \Cref{prop:e-relations} to our situation, $J_w$ is generated by the following elementary symmetric polynomials\footnote{Each of these relations can alternatively be obtained from a simple Chern class analysis on the tautological quotient line bundles on $\fl_n$.}:
    \begin{itemize}
        \item $e_m(x_1,\dots, x_r)$ for $r<i$, $m = 1,\dots, r$,
        \item $e_m(x_1,\dots, x_{w^{-1}(i)})$ for $m=i-w^{-1}(i)+1,\dots, w^{-1}(i)$,
        \item $e_m(x_{w^{-1}(j)},\dots, x_n)$ for $m = j-w^{-1}(j)+1,\dots, n-w^{-1}(j)-1$,
        \item $e_m(x_{r+1},\dots, x_n)$ for $r>j$, $m=1,\dots, n-r$.
    \end{itemize}
    The first and last set of generators can be simplified to generators $x_1,\dots, x_{i-1}$ and $x_{j+1},\dots, x_n$. Modding out by these variables, we see that the second family of generators then simplifies to $e_{i-w^{-1}(i)+1}(x_i,\dots, x_{w^{-1}(i)}) = x_i\cdots x_{w^{-1}(i)}$. Furthermore, the third set of generators reduces to $e_{j-w^{-1}(j)+1}(x_{w^{-1}(j)},\dots, x_j) = x_{w^{-1}(j)}\cdots x_j$. 
    Thus, $J_w$ is generated by the squarefree monomials stated above.
\end{proof}





\begin{lemma}[\Cref{cond:A2} implies \Cref{cond:A3}]\label{lem:2implies3}
\label{lem: A2 implies A3}
    Suppose $w\in S_n(\calP)$, then $X_w$ is defined by the inclusions in \Cref{cond:A3}.
\end{lemma}

\begin{proof}
    If $w\in S_n(\calP)$, then $w$ also avoids \(\{4231, 35142, 42513,  351624\}\). Thus, $X_w$ is defined by inclusions in the sense of Gasharov--Reiner \cite{GasharovReiner}. In fact, the inclusion conditions for $w\in S_n(\calP)$ take the following precise form.

    If $w=\id$, then $X_\id$ is a single point consisting of the unique flag defined by $V_b=E_a$ for all $i$, where $E_\bullet$ is the standard flag. Otherwise, if $w\neq \id$ then by \Cref{lem:AvoidEquivDecrease}, the letters $i+1,i+2,\dots, j-1$ are in decreasing order in $w$, where $i$ and $j$ are the smallest and largest non-fixed-points of $w$, respectively. It is not hard to check that if $w^{-1}(i)\neq j$ and $w^{-1}(j)\neq i$, then the essential set $\mathrm{Ess}(w)$ of $w$ is
    \begin{gather*}
    \{(2,1),(3,2),\dots, (i,i-1)\}\,\cup\, \{(i+1,w^{-1}(i))\}\\ 
    \cup\,  \{(j,w^{-1}(j)-1)\}\, \cup \,\{(j+1,j),(j+2,j+1),\dots, (n-1,n)\}.
    \end{gather*}
    See \Cref{ex:Essential} for the case of $w=1657423$ where $i=2$ and $j=7$. 
    
    Each of these sets immediately translates into the inclusion conditions in \Cref{cond:A3}. For example, for $k<i$, $(k+1,k)\in \mathrm{Ess}(w)$ and $r_w(k,k) = k$. Thus, a defining condition of $X_w$ is $\dim_{\bC}(V_k\cap E_k)\geq k$, forcing $V_k = E_k$. For $(i+1,w^{-1}(i))$, $r_w(i,w^{-1}(i)) = i$, so $\dim_\bC(V_{w^{-1}(i)}\cap E_i)\geq i$, forcing the containment $E_i\subseteq V_{w^{-1}(i)}$. The other pairs in the essential set similarly translate to the other two sets of containment conditions in \Cref{cond:A3}. 
    
    In the case $w^{-1}(j)=i$, the pair $(j,w^{-1}(j)-1)$ is not in the essential set. This is simply because the corresponding condition $\dim_{\bC}(V_{i-1}\cap E_{j-1})\geq i-1$ translates to $V_{i-1}\subseteq E_{j-1}$, which is redundant from the condition $V_{i-1} = E_{i-1}$. Similarly, if $w^{-1}(i) = j$ then the containment condition becomes $E_i\subseteq V_j$, which is redundant from $V_j = E_j$.
\end{proof}

We record here a consequence of combining the previous two lemmas. Recall that $\pi:S\to R_n$ is the usual projection map.

\begin{lemma}
    \label{lem:defined by inclusions}
    If $w\in S_n(\calP)$, then $\pi^{-1}(J_w) = I_{\Delta_w}+I_n$.
\end{lemma}
\begin{proof}
    Suppose $w\in S_n(\calP)$. Combining \Cref{lem: A2 implies A3} and \Cref{lem: A3 implies A1}, we have that $J_w$ is generated by the squarefree monomials listed in \Cref{lem: A3 implies A1}. 
    %By \Cref{lem:2implies3}, if $w\in S_n(\calP)$ then it is defined by the inclusions in \Cref{cond:A2}.
    %As the minimal $v$ such that $v\not\leq w$ each come from a left ascent of $w$, there are at most $n-j+i+1$ such $v$.
    %Then by \cite[Example 4.9]{ReinerWooYong}, the generators of the ideal $J_w$ are given by $\{e_m(x_1,\dots,x_r):m=r-s+1,\dots,r\}$ for each defining Schubert condition $E_s\subseteq V_r $ and $\{e_m(x_{r+1},\dots,x_n):m=s-r+1,\dots,n-r\}$ for each defining Schubert condition $V_r \subseteq E_s$.
    %In particular, since $e_1(x_1,\dots,x_k)\in J_w$, then $x_k\in J_w$ for $k\in[1,i-1]$.
    %Similarly, since $e_1(x_{k+1},\dots,x_n)\in J_w$, then $x_k\in J_w$ for $k\in[j+1,n]$.
    %Those inclusions generate all of the other generators coming from the equalities as their successive differences show that the variables $x_k$ for $k\in [1,i-1]\cup [j+1,n]$ are in the ideal.
    %Now we only have the inclusions $E_i\subseteq V_p$ and $V_{q-1}\subseteq E_{j-1}$ left.
    %They contribute generators $\{e_m(x_1,\dots,x_p):m=p-i+1,\dots,p\}$ and $\{e_m(x_{q},\dots,x_n):m=j-q+1,\dots,n-q+1\}$.
    %Due to the variable inclusions, these are equivalent to $\{e_m(x_i,\dots,x_p):m=p-i+1,\dots,p\}$ and $\{e_m(x_{q},\dots,x_{j}):m=j-q+1,\dots,n-q+1\}$.
    %However, $x_i$ through $x_p$ is only $p-i+1$ variables so, in fact, there is only one new generator from that inclusion, $e_{p-i+1}(x_i,\dots,x_p)$.
    %Similarly, the other inclusion only contributes the new generator $e_{j-q+1}(x_{q},\dots,x_{j})$.
    %
    %Both of these new generators (as well as the variable generators above) are squarefree monomials.
    By inspection, these are exactly the monomials corresponding to the minimal nonfaces of $\Delta_w$, and so the result follows.
\end{proof}

Before continuing to complete the proofs of \Cref{thm:main-A} and \Cref{thm:main-B}, we count the exact number of minimal generators of $J_v$ for $v\in \sn(\calP)$.
\begin{proposition}
\label{prop: min gens of Iv}
    Let $v\in \sn(\calP)$ with $v\neq \mathrm{id}$, and let $i=i(w)$, $j=j(w)$, $p=p(w)$, and $q=q(w)$. 
    If $p \neq j$ or  $q\neq i$, then the minimal generating set of $I_{\Delta_v}$ given by the minimal nonfaces of $\Delta_v$ descends to a minimal generating set of $J_v$ in $R_n$. 
    If $p = j$ and $q= i$, then the images of the variables $x_1,\dots, x_{i-1}, x_{j+1},\dots, x_n$ in $R_n$  
    %the variable minimal generators of $I_{\Delta_v}$ 
    are a minimal generating set of $J_v$.
\end{proposition}
\begin{proof}
    By \Cref{lem:defined by inclusions}, the generating set of $I_{\Delta_v}$ descends to a generating set of $J_v$ so it suffices to show that they form a minimal generating set. 
    As there must be at least two variables that are not generators, $e_1$ cannot make any of the variable generators redundant.
    Thus, without loss of generality, we may assume that $i=1$ and $j=n$.
    
    If $p=j$ and $q=i$, then the only generator of $I_{\Delta_v}$ is $e_n$ so $J_v = (0)$, and the result holds in this case.

    If $p\neq j$ or $q\neq i$, it suffices to show that the images of the generators of $I_{\Delta_v}$ are nonredundant.
    Recall that any $w\in \calP$ is determined by the tuple $[i_w, j_w, p_w,q_w]$.
    Then $x_{[i,p]}$, $x_{[q,j]}$, or both is a minimal generating set.
    
    If $p\neq j$, then $x_{[i,p]}$ is a member of a minimal generating set of $I_{\Delta_v}$.
    If it was not a member of a minimal generating set of $J_v$, then $J_v \subseteq J_w$ where $w\in \calP$ is  determined by the tuple $[i_w, j_w, p_w,q_w] = [i_v, j_v, j_v,q_v]$.
    As $v\leq w$, $J_w\subset J_v$ as well by \Cref{prop:Inclusion} and \Cref{lem:defined by inclusions}, so $J_v=J_w$, which is impossible as $v\neq w$.
    Thus, $x_{[i,p]}$ is a member of a minimal generating set of $J_v$ (if it was a member of a minimal generating set of $I_{\Delta_v}$).

    A symmetric argument shows that $x_{[q,j]}$ is a member of a minimal generating set of $J_v$ (if it was a member of a minimal generating set of $I_{\Delta_v}$).
\end{proof}

\begin{remark}
\label{rem: number of min gens of Iv}
    It follows from the previous result that $J_w$ has a set of $n-j+i+1-\delta_{jp}-\delta_{iq}$ minimal generators in $R_n$. 
    For example, if $w = 3412,$ then $I_{\Delta_w}$ has two generators: $x_1x_2x_3$ and $x_2x_3x_4,$ since $i = 1, p = 3, q = 2,$ and $j = 4$ as in \Cref{def:Deltaw}. Thus, $n - j + i + 1 - \delta_{jp} - \delta_{iq} = n - m + 2 = 2,$ where $m := j - i + 1 = \#\supp(\Delta_w).$ One can also see that $J_w$ has a set of $n+\delta_{q<p}$ minimal generators in $R_n$ as follows.

    Note that any variable that is a generator of $J_w$ makes exactly one of the generators of $I_n$ redundant, so without loss of generality, we may assume that $J_v$ has no variable generators.
    If $J_w$ has a single generator, it is either $e_n$ or makes $e_n$ redundant (and no other generators of $I_n$ as each contains a monomial divisible by the generator) so the result holds in this case.  
    
    If $J_w$ has two minimal generators with $p<q$, then they make $e_{n-1}$ and $e_n$ redundant as any monomial missing at most one variable must be divisible by one or the other.
    No other generators of $I_n$ are made redundant as again they all contain a monomial not divisible by either generator (any monomial excluding one variable from each generator).
    Thus, the result holds in this case as well.
    
    If $J_w$ has two minimal generators with $p>q$, then they $e_n$ redundant, but do not make any other generator of $I_n$ redundant as again they all contain a monomial not divisible by either generator (any monomial excluding a variable from the overlapping variables). 
\end{remark}

We can now prove the first implication of \Cref{thm:main-B}.

\begin{proof}[Proof that \Cref{cond:B2} implies \Cref{cond:B1}]
   Let $\mathcal{W}\subseteq S_n(\mathcal{P})$ such that 
   \[\bigcap_{w\in\mathcal{W}} (I_{\Delta_w}+I_n) = \left(\bigcap_{w\in \mathcal{W}} I_{\Delta_w}\right)+I_n.\]
   Letting $\Delta_\mathcal{W} = \bigcup_{w\in\mathcal{W}}\Delta_w$, we then have 
   \begin{align*}
   \pi^{-1} \left( J_{\mathcal{W}} \right) 
   &= \pi^{-1}\left(\bigcap_{w\in\mathcal{W}} J_w\right)\\
   &= \bigcap_{w\in\mathcal{W}} \pi^{-1}(J_{w})\\
   &= \bigcap_{w\in\mathcal{W}} \left(I_{\Delta_w}+I_n\right) \text{ by \Cref{lem:defined by inclusions}}\\
   &= \left(\bigcap_{w\in\mathcal{W}} I_{\Delta_w}\right)+I_n \\
   & = I_{\Delta_\mathcal{W}}+I_n.
   \end{align*}
Applying the surjective map $\pi$ then gives 
\[J_\mathcal{W} = \pi\left(\pi^{-1}(J_{\mathcal{W}})\right) = \pi\left( I_{\Delta_\mathcal{W}}+I_n \right) = \overline{I_{\Delta_\mathcal{W}}},\]
so $J_\mathcal{W}$ is generated by squarefree monomials as desired.
\end{proof}

%------------------------------------------------------------
\subsection{Squarefree generation implies ideal equality}
%------------------------------------------------------------
We prove this implication of \Cref{thm:main-B} in two steps.
We first show the following.

\begin{proposition}
\label{prop:1-implies-avoidance}
If \Cref{cond:B1} holds, then $\mathcal W\subseteq S_n(\mathcal{P})$.
\end{proposition}

Following the proof of this proposition, we give that proof that \Cref{cond:B1} implies all of \Cref{cond:B2}.

\begin{proof}
    Let $\mathcal{W}\subseteq S_n$ be a set of permutations for which $J_{\mathcal{W}} = \bigcap_{w\in\mathcal W} J_w$ is generated by squarefree monomials as an ideal of $R_n$.
    Let $w\in\mathcal{W}$.
    Suppose for contradiction that $w\not\in S_n(\mathcal{P})$.

    Set \(\ell=\ell(w)\). 
    Recall to each permuation $w\in S_n$, we associate the set $A(w)=\mathrm{Des}_L(w)^c\cap[i+1,j-2]$.
    Using this, define a linear functional on \((R_n)_{\ell+1}\) by 
    \[  \Lambda_w:=\sum_{q\in A(w)}\lambda_{s_qw}. \]
    By \Cref{lem:AvoidEquivDecrease}, there exists some $q\in A(w)$ so this sum is non-empty.
    Fix a $q\in A(w)$.
    By definition of the $\lambda_*$, we have
    \[ \Lambda_w(\mathfrak{S}_{s_qw})=\sum_{q'\in A(w)}\lambda_{s_q'w}(\mathfrak{S}_{s_qw}) = \lambda_{s_qw}(\mathfrak{S}_{s_qw})=1. \]
    Thus, in order to find a contradiction, it suffices to show that we can also prove $\Lambda_w(\mathfrak{S}_{s_qw}) = 0$.


    In order to show that, we first need to show that $\mathfrak{S}_{s_qw}\in \mathrm{SqfMon}(J_w)$, which we do in this paragraph, and then show that $\Lambda_\mathcal{W}$ evaluates to $0$ on that entire ideal.
    If $s_qw\leq v\in\mathcal{W}$, then $w< s_qw\leq v$ which contradicts that $\mathcal{W}$ is an antichain (or the reflexivity of the strong Bruhat order).
    Thus, $s_qw\not\leq v$ for all $v\in \mathcal{W}$, so $\mathfrak{S}_{s_qw}\in J_{\mathcal{W}}$.
    As $J_{\mathcal{W}}\subseteq J_w$ and $J_{\mathcal{W}}$ is generated by squarefree monomials, $J_{\mathcal{W}}\subseteq \mathrm{SqfMon}(J_w)$.
    Thus, $\mathfrak{S}_{s_qw}\in \mathrm{SqfMon}(J_w)$.
    Thus, in order to find a contradiction by showing $\Lambda_w(\mathfrak{S}_{s_qw}) = 0$, it suffices to show that for all $f\in \mathrm{SqfMon}(J_w)$,  $\Lambda_{w}(f) = 0$.
    
    As $\mathrm{SqfMon}(J_w)$ is generated by squarefree monomials, $\mathrm{SqfMon}(J_w)_{\ell+1}$ is generated as a vector space by products of the form $hx_F$ where $x_F\in \left(\mathrm{SqfMon}(J_w)\right)_d \subseteq (J_w)_d$  and $h\in (R_n)_{\ell+1-d}$.
    Thus, to show that for all $f\in \mathrm{SqfMon}(J_w)$,  $\Lambda_{w}(f) = 0$, it suffices to show that $\Lambda_{w}(hx_F) = 0$ for all such pairs.

    By definition of $\Lambda_w$,  Remark \ref{rem:lin-func-length}, and \cite[Formula (2.7)]{Macdonald}, and Definition \ref{def:relative-diff-oper}, respectively, we have

    \begin{align}
    \label{eq: functional string of identities}
    \Lambda_w(hx_F)
    & = \sum_{q\in A(w)}\!\!\lambda_{s_qw}(hx_F) \\
    & = \sum_{q\in A(w)}\!\!\partial_{s_qw}(hx_F) \\
    & = \sum_{q\in A(w)}\!\!\partial_q\partial_w(hx_F) \\ 
    & = \sum_{q\in A(w)}\!\!\partial_q\sum_{r\leq w}r\!\left(\partial_{w/r}h\right)\partial_r(x_F)    
    \end{align}
    
    Now, $\deg\!\left(\partial_r(x_F)\right) = d-\ell(r)$ and \[ \deg\!\left(r(\partial_{w/r}h)\right) = \deg\!\left(\partial_{w/r}h\right) = \deg(h)-\bigl(\ell-\ell(r)\bigr) = \ell(r)+1-d.\] 
    In order for $r$ to possibly contribute a nonzero term, those must both be nonnegative so $d-1\leq \ell(r) \leq d$.

    When $\ell(r) =d$, $\partial_r(x_F) = \lambda_r(x_F)$ by Remark \ref{rem:lin-func-length}. 
    Since $x_F\in J_w$ and $J_w$ is spanned as a vector space by Schubert polynomials $\mathfrak{S}_{u}$ for $u\not\leq w$, then $\lambda_r(x_F) = 0$ for all $r\leq w$. 
    %For all $r\leq w$, $\lambda_r(x_F) = 0$ as $x_F\in J_w$ so in the basis of Schubert polynomials its coefficients are $0$ for all of the Schubert polynomials $\mathfrak{S}_r$ where $r\leq w$.
    Thus, the terms in \eqref{eq: functional string of identities} with $\ell(r) =d$ contribute $0$, so \[\Lambda_w(hx_F) = \sum_{q\in A(w)}\!\!\partial_q\!\sum_{\substack{r\leq w\\\ell(r)=d-1}}\!\!\!r\left(\partial_{w/r}h\right)\partial_r(x_F).\]
    
    Then $\partial_{w/r}h$ has degree $\ell(r)-\ell+\ell+1-d = 0$, so it is a scalar, and hence is fixed by $r$, so the sums can be reordered as 
    \[
    \Lambda_w(hx_F) = \sum_{\substack{r\leq w\\\ell(r)=d-1}}\!\!\!r\left(\partial_{w/r}h\right)\sum_{q\in A(w)}\!\!\partial_q\partial_r(x_F)=0.
    \]
    Thus, it in order to show $\Lambda_{w}(hx_F) = 0$ for all relevant pairs $h$ and $x_F$, it  suffices to show $\sum_{q\in A(w)}\partial_q\partial_r(x_F)=0$ for all $\ell(r) = d-1$ and $r\leq w$. 
    If \(q'\in A(w)^C \subseteq [i+1,j-2]\), then \(s_{q'}w<w\), then \(\partial_{q'}\partial_w=0\) by \cite[Equation 2.7]{Macdonald} so it further suffices to show \[\displaystyle\sum_{q\in [i+1,j-2]}\partial_q\partial_r(x_F)=0\] when $\ell(r) = d-1$ and $r\leq w$ (and $x_F\in J_w$ is a squarefree monomial of degree $d$). 
    
    
    Fix $r\leq w$ with $\ell(r) = d-1$ and $x_F \in (J_w)_d$ a squarefree monomial.
    Writing $r$ as a reduced word and repeatedly applying Remark \ref{rem:sqfr-mon-div-diff} gives that $\partial_r(x_F) = c_F x_a$ for some $c_F\in \{0,\pm 1\}$ and some $a\in [n]$.
    If $c_F = 0$, we have reached our conclusion so assume without loss of generality that $c_F \neq 0$.

    If $s_{a-1}$ or $s_{a}$ are in any reduced word expression for $r$, then  $c_F = 0$ as either of those would have eliminated that variable $x_a$ from $x_{F}$ either by dividing it out or by zeroing it out.
    Thus, we have 
    \begin{equation}
    \label{eq:two-terms}
     \partial_q\partial_r(x_F)  =
     \begin{cases}
     -c_F,&q=a-1,\\
     c_F,&q=a,\\
     0,&\text{otherwise}.
     \end{cases}
     \end{equation}
In other words, the sum simplifies to 
\[\displaystyle\sum_{q\in [i+1,j-2]}\partial_q\partial_r(x_F)=\displaystyle\sum_{q\in [a-1,a]\cap [i+1,j-2]}\partial_q\partial_r(x_F).\]
If $[a-1,a]\cap [i+1,j-2] = \varnothing$, then the sum is trivially zero as desired.
If $[a-1,a]\cap [i+1,j-2] = [a-1,a]$, then $1$ and $-1$ each appear once in the sum so it is zero as desired.

In order to show that the sum is $0$, it now suffices to show that if $|[a-1,a]\cap [i+1,j-2]| = 1$, then there is a contradiction.
Suppose for contradiction that $|[a-1,a]\cap [i+1,j-2]| = 1$.
Then either $a=i+1$ or $a-1 = j-2$; in the second case, $a=j-1$.
Set $b = i = a-1$ or $b =j-1=a$, respectively.
Let $ \mathrm{supp}(v)$ denote the set of $i$ such that $s_i$ appears in some/every reduced word giving $v$.
Since \(r\leq w\),
\[\mathrm{supp}(r)\subseteq \mathrm{supp}(w)\subseteq\{i,\ldots,j-1\}. \]
As $c_F\neq 0$, neither $a$ nor $a-1$ is in $\mathrm{supp}(r)$.
So in our two cases, $i,i+1\not\in \mathrm{supp}(r)$ and $j-2,j-1\not\in \mathrm{supp}(r)$, respectively.
In either case, $s_b$ commutes with every reduced word for $r$.

By commutativity and the fact that $b\not\in\mathrm{supp}(r)$, appending $s_b$ to a reduced word for $r$ gives a reduced word for $s_br$.
As $b\in\{i,j-1\}$ and $w$ fixes every $k\in[i,j]$ but neither $i$ nor $j,$ the permutation $w$ does not fix $[1,b],$ so $s_b$ occurs in every reduced word of $w.$ Fix a reduced word for $w.$ Since $r\leq w,$ by the subword property \cite[Corollary 2.2.3]{BjornerBrenti} some subword of it is a reduced word for $r,$ and as $b\not\in\supp(r),$ this subword uses no occurrence of $s_b.$ If we include one occurrence of $s_b$ into the subword, then we obtain a subword of the reduced word for $w,$ whose product is $s_br,$ since $s_b$ commutes with every letter of $r,$ and which is reduced by the previous sentence. Hence, $s_b\leq w.$


Recall that the Schubert polynomials form a basis for $R_n$ and the Schubert polynomials for $v\not\leq w$ form a basis of $J_w$.
Because \(x_F\in J_w\), it can be written in terms of that subbasis (so its coefficients on the other Schubert polynomials are $0$).
Thus,  \(\partial_{s_br}(x_F)=\lambda_{s_br}(x_F)=0\).  
In contrast, by \eqref{eq:two-terms}
\[  
\partial_{s_br}(x_F)=\partial_b\partial_r(x_F) =\pm c_F. 
\] 
This contradiction rules out the case of exactly one index lying in the set, so the sum is zero as desired.
As the sum is $0$, $\Lambda_{w}(\mathrm{SqfMon}(I))=0$.
As $\mathfrak{S}_{s_qw} \in \mathrm{SqfMon}(J_w)$, we have $\Lambda_{w}(\mathfrak{S}_{s_qw}) = 0 = 1$, and we have the contradiction we needed.
\end{proof}


We need one more lemma before proving the second implication 

\begin{proposition}
\label{prop:lift}
If \(w\in S_n(\mathcal{P})\) and \(x_F\) is squarefree, then
\[ 
\overline{x_F}\in J_w \quad\Longleftrightarrow\quad  x_F\in I_{\Delta_w}, 
\]
where $\overline{x_F}$ is the image of $x_F$ under the projection $\pi$.
\end{proposition}

\begin{proof}
If $w =\mathrm{id}$, then $I_{\Delta_w}$ and $J_w$ are each $(x_1,\dots,x_n)$ in their respective rings, so the result is immediate as they both contain all $x_F$ and do not contain $1$.
For the rest of the proof, let $w\neq \mathrm{id}$.


The reverse implication is also immediate since $\pi^{-1}(J_w)=I_{\Delta_w}+I_n$ by \Cref{lem:defined by inclusions}.  

For the forward implication, we work by contrapositive, so suppose that \(x_F\notin I_{\Delta_w}\).  
By definition of $I_{\Delta_w}$, we have
\[ F\subsetneq[i,j],\qquad [i,p]\nsubseteq F,\qquad [q,j]\nsubseteq F.\]
Without loss of generality, mod out by the variables $x_a$ for  \(a\in [i,j]^c\subset[n]\).  
By \Cref{lem:quotient-and-annihilators}, Part (1), the quotient is \(R_m\), where \(m=j-i+1\), and the two remaining generators\footnote{After relabeling $x_{k+i-1} = y_{k}$} of \(J_w\) are
\[  
y_G=y_1\cdots y_{p-i+1}\qquad\text{and}\qquad y_H=y_{q-i+1}\cdots y_m.
\]
Here \(p>i\) and \(q<j\), so both degrees are at least two.  
Further, as $p\neq q$, their overlap does not have size 1.
In terms of monomials, the three relations above become  
\[
y_F\mid y_{[1,m]},\qquad y_G\nmid y_F,\qquad y_H\nmid y_F,\qquad |G\cap H|\neq 1.
\]

By \Cref{lem:contraction}, $\overline{y_F}\not\in (\overline{y_G},\overline{y_H}) = J_w \subseteq R_m$.
Lifting back from the quotient, the contrapositive follows.

\iffalse
\tim{Shouldn't the rest follow faster from Lemma 3.3?}
If $y_G = y_H$, then $G = H$; since $G$ has left endpoint $1$ and $H$ has right endpoint $m$, this forces $G = H = [1,m]$, and hence
$\overline{y_G} = \overline{y_H} = \overline{y_1 \cdots y_m} = \overline{e_m} = 0$ in $R_m$, while $\overline{y_F}$ is a proper squarefree monomial and so is nonzero by \Cref{lem:ann-squarefree-proofonly}. Thus, $\overline{y_F}\not\in (\overline{y_G},\overline{y_H}) = J_w \subseteq R_m$.

If \(y_G\neq y_H\), $G\cap H=\varnothing$ when \(p<q\) and $|G\cap H|\geq 2$ when \(p>q\).  
Note, the case \(p=q\) is impossible, as \(w^{-1}(i)\neq w^{-1}(j)\). 
As $y_G\nmid y_F$ and $y_H\nmid y_F$,  \(\overline{y_F}\not\in(\overline{y_G},\overline{y_H}) = J_w\) by Lemma \ref{lem:contraction}.  
This proves the contrapositive.
\fi
\end{proof}

Given \Cref{prop:1-implies-avoidance} and \Cref{prop:lift}, we can now prove the second implication of the main result. From here on, if $I$ is an ideal in $S$, we will use $\overline{I}$ to denote the corresponding ideal $\pi(I)$ in $R_n$.


\begin{proof}[Proof that \Cref{cond:B1} and $\mathcal{W}\subseteq S_n(\mathcal{P})$ implies \Cref{cond:B2}.]
Assume that $J_{\mathcal{W}}$ is generated by squarefree monomials, i.e., \(J_\mathcal{W} = \{\overline{m_1},\dots,\overline{m_r}\}\subseteq R_n\) where \(\{m_1,\dots,m_r\}\subset \mathbb{C}[x_1,\dots,x_n]\) are squarefree monomials, and $\mathcal{W}\subseteq S_n(\mathcal{P})$.


By Proposition \ref{prop:lift}, \(m_a\in I_{\Delta_w}\) for all \(w\in\mathcal{W}\) and \(a\in[r]\), so \(m_a \in  \bigcap_{w\in \mathcal{W}} I_{\Delta_w}\).
Thus, \(\overline{m_a}\in \overline{\left(\bigcap_{w\in \mathcal{W}}  I_{\Delta_w}\right)}\).
As all of the generators of $J_{\mathcal{W}}$ are contained in $\overline{\left(\bigcap_{w\in \mathcal{W}}  I_{\Delta_w}\right)}$, we have the inclusion $J_{\mathcal{W}}\subseteq \overline{\left(\bigcap_{w\in \mathcal{W}}  I_{\Delta_w}\right)}$.
Using that inclusion and that \(\overline{I_{\Delta_w}}=J_w\) by \Cref{lem:defined by inclusions}, we have the following string of inclusions:
\[ J_{\mathcal W} \subseteq \overline{\left(\bigcap_{w\in\mathcal W}I_{\Delta_w}\right)} \subseteq \bigcap_{w\in\mathcal W}\overline{I_{\Delta_w} }=\bigcap_{w\in\mathcal W}J_{w} = J_{\mathcal{W}}.\]
Thus, all of these are equal.

Finally, as $\pi^{-1}(J_w) = I_{\Delta_w}+I_n$ by \Cref{lem:defined by inclusions}, we have
\[
\begin{aligned}
\bigcap_{w\in\mathcal W}(I_{\Delta_w}+I_n)
&=\bigcap_{w\in\mathcal W}\pi^{-1}(J_w) =\pi^{-1}\!\left(\bigcap_{w\in\mathcal W}J_w\right)
=\pi^{-1}(J_{\mathcal W})\\
&=\pi^{-1}\left(\overline{\bigcap_{w\in\mathcal W}I_{\Delta_w}}\right)
=\left(\bigcap_{w\in\mathcal W}I_{\Delta_w}\right)+I_n.
\end{aligned}
\] 

Thus, \Cref{cond:B2} holds as desired.
\end{proof}







%----------------------------------------------
\subsection{Ideal equality implies prunable}
%----------------------------------------------
In this section, we prove the third implication of \Cref{thm:main-B} by induction on the length of $\mathcal{W}$.
As both \Cref{cond:B1,cond:B2} require the set to be contained in $S_n(\calP)$, and all of their other parts are trivial for singletons, that base case is trivial.
However, in the proof we need the additional base case of when the set has size 2.

\subsubsection{Ideal equality implies prunable for pairs.}
For pairs, \Cref{cond:B3}(b) is trivial, so it suffices to show \Cref{cond:B3}(a).

We first prove a more general lemma about intersections of ideals of simplicial complexes plus $I_n$ and the dimension of the intersection of the complexes.
When applied in our setting, this gives the dimension aspect of (3)(a). 
\begin{lemma}
\label{lem:dimension-complete}
If \(\Delta_{1},\Delta_{2}\) are simplicial complexes such that $(I_{\Delta_{1}}+I_n)\cap(I_{\Delta_{2}}+I_n)  =(I_{\Delta_{1}}\cap I_{\Delta_{2}}) +I_n$, then $ \dim(\Delta_{1}\cap\Delta_{2})  =\min\{\dim\Delta_{1},\dim\Delta_{2}\}$.
\end{lemma}


\begin{proof}
  Start with the exact sequence
    \[
    0\longrightarrow S/I_{\Delta_{1}\cup\Delta_{2}} \longrightarrow S/I_{\Delta_{1}}\oplus S/I_{\Delta_{2}} \longrightarrow S/I_{\Delta_{1}\cap\Delta_{2}} \longrightarrow0 
    \]
where $S$ is the polynomial ring.
Assume for contradiction that 
\[
d=\dim(\Delta_{1}\cap\Delta_{2}) <\min\{\dim\Delta_{1},\dim\Delta_{2}\}.
\]
Since \(\Delta_{1}\cap\Delta_{2}\) has no face of dimension \(d+1\), the element \(e_{d+2}\) vanishes in \(S/I_{\Delta_{1}\cap\Delta_{2}}\). 
Therefore, \((e_{d+2},0)\) is in the kernel of the surjective map.
By exactness, there exists an element \(\xi\in S/I_{\Delta_{1}\cup\Delta_{2}}\) mapping to \((e_{d+2},0)\).
Note, \((e_{d+2},0)\) is nonzero as $\Delta_{1}$ has a face of dimension at least $d+1$.

By the assumption on the ideal, we can tensor by $R_n$ to get the sequence
\[0\longrightarrow R_n/I_{\Delta_{1}\cup\Delta_{2}} \longrightarrow R_n/I_{\Delta_{1}}\oplus R_n/I_{\Delta_{2}} \longrightarrow R_n/I_{\Delta_{1}\cap\Delta_{2}} \longrightarrow0 \]

In the quotient sequence, we have \(\overline{(e_{d+2},0)} = \overline{(0,0)}\), so by injectivity \(\overline{\xi} =\overline{0}\).
Thus, \(\xi\in \overline{I_n}\subseteq S/I_{\Delta_{1}\cup\Delta_{2}}\), so we can write
\[ \xi=\sum_{j=1}^{d+1}e_jh_j+c e_{d+2} \]
where $c$ is a scalar.
Taking the two coordinate images in the original short exact sequence gives
\[ e_{d+2}=\sum_{j=1}^{d+1}e_jh_j+c e_{d+2} \quad\text{in }S/I_{\Delta_{1}}\qquad\text{and}\qquad 0=\sum_{j=1}^{d+1}e_jh_j+c e_{d+2} \quad\text{in }S/I_{\Delta_{2}}. \]


As $\Delta_{1}$ has a $d+1$-dimensional, $e_{d+2}$ is not in the ideal generated by the smaller elementary symmetric polynomials\footnote{Quotient out by the variables outside the the $d+1$ dimensional face and this becomes the standard statement on the ring $R_m$}.
Thus, by the 1st equality, \(c\neq 0\in S/I_{\Delta_1}\).
Applying the same logic to the 2nd equality, gives \(c=0\in S/I_{\Delta_2}\).
As $c$ was a scalar, those two statements are contradictory, and so the result follows.
\end{proof}

Next, we prove the purity aspect of (3)(a).
\begin{lemma}
\label{lem:crossing-class-complete}
Let \(\{u,v\}\subseteq S_n(P)\) be an antichain for which $(I_{\Delta_u}\cap I_{\Delta_v}) +I_n = (I_{\Delta_u}+I_n)\cap (I_{\Delta_v}+I_n)$.
Then \(\Delta_u\cap\Delta_v\) is pure.
\end{lemma}


\begin{proof}
    We work by contradiction so assume that \(\Delta_u\cap\Delta_v\) is not pure.
    As the intersection of doubly pruned complexes is a doubly pruned complex, \(\Delta_u\cap\Delta_v\) is doubly pruned, i.e. its minimal nonfaces are vertices outside of some interval $[i,j]$ and, possibly, two interval nonfaces of the form $[i,p]$ and $[q,j]$.
    Since it is not pure, $q\leq p$.
    As $\{u,v\}$ is an antichain, assume without loss of generality that $[i,p]$ is a nonface of $\Delta_u$ and a face of $\Delta_v$ and $[q,j]$ is a nonface of $\Delta_v$ and a face of $\Delta_u$.
    
    Define the permutation $w$ as follows: 
    \[
    w :=(1,\dots,i-1,i+1,\dots,q-1,j,q,\dots,p,i,p+1,\dots,j-1,j+1,\dots,n).
    \]
    We then have the rank inequality 
    \[
    r_w(p,i) = i-1 <i = r_u(p,i)
    \]
    so $w\not\leq u$.
    Similarly, we have 
    \[
    r_w(q-1,j-1)=q-2<q-1 = r_v(q-1,j-1)
    \]
    so $w\not\leq v$.
    Since $w\not\leq u,v$, it follows that $\mathfrak{S}_w\in J_u\cap J_v$ in $R_n$.
    
    By \Cref{lem:defined by inclusions}, $\pi^{-1}(J_u) = I_{\Delta_u}+I_n$ and $\pi^{-1}(J_v) = I_{\Delta_v}+I_n,$ so by our hypothesis,
    \[
    \pi^{-1}(J_u\cap J_v) \subseteq \pi^{-1}(J_u)\cap \pi^{-1}(J_v)
    = (I_{\Delta_u}+I_n)\cap(I_{\Delta_v}+I_n) 
    = (I_{\Delta_u}+I_{\Delta_v}) + I_n
    \]
    and in particular any lift of $\fraks_w$ to $S$ lies in $(I_{\Delta_u}\cap I_{\Delta_v})+I_n.$ 
    It now suffices to show that $\fraks_w\not\in\pi\left((I_{\Delta_u}\cap I_{\Delta_v})+I_n\right),$ i.e., that no lift of $\fraks_w$ lies in $(I_{\Delta_u}\cap I_{\Delta_v})+I_n.$
    Consider the permutation
    \[
    t := (1,\dots,i-1, 
    j-1,\dots,j-q+i+1,
    j, 
    j-q+i,\dots,j-p+i, 
    i, j-p+i-1,
    \dots,i+1,
    j+1,\ldots,n).
    \]
    Notice that $t\in S_n(\calP)$ is the unique strong Bruhat maximal element among any of the permutations $\sigma$ with $\sigma^{-1}(i)=p+1$, $\sigma^{-1}(j)=q-1$, such that $\sigma(k) = k$ for all $k\not\in[i,j]$. By \Cref{prop: Ehresmann's criterion}, placing the remaining values $i+1,\ldots,j-1$ in decreasing order in the remaining positions minimizes every entry of the rank matrix among such permutations, and $t\in\sn(\calP)$ by \Cref{lem:AvoidEquivDecrease}. 
    Since $w$ satisfies those conditions, $w\leq t$, and so $\fraks_w\not\in J_t,$ i.e., no lift of $\fraks_w$ lies in $\pi^{-1}(J_t) = I_{\Delta_t}+I_n.$ Thus, it suffices to show that $(I_{\Delta_u}+I_{\Delta_v})+I_n\subseteq I_{\Delta_t}+I_n,$ and for this it suffices to show that $I_{\Delta_u}+I_{\Delta_v}\subseteq I_{\Delta_t}.$
    
    As the vertex nonfaces of $u$ and $v$ are clearly both contained in the vertex nonfaces of $t$, we only need to show that the 4 least common multiples of nonsingleton generators in $I_{\Delta_u}$ and $I_{\Delta_v}$ are in $I_{\Delta_t}$ (recall that $I_{\Delta_u}\cap I_{\Delta_v}$ is generated by the least common multiple of pairs of generators, one from each ideal).

    The left end interval of $v$ either contains a vertex nonface of $t$ or contains the nonface $[i,p+1]$ so any least common multiple using that face is in $I_{\Delta_t}.$
    Similarly, the right end interval of $u$ either contains vertex nonface of $t$ or contains the nonface $[q-1,j]$ so any least common multiple using that face is in $I_{\Delta_t}$.
    Finally, the least common multiple of the generators coming from nonfaces $[i,p]$ and $[q,j]$ is the monomial $x_{[i,j]}$ which is clearly in $I_{\Delta_t}$.
    Thus, we have the needed inclusion, and the result follows.
\end{proof}

\begin{lemma}
\label{lem:2 implies 3 for pairs}
    If $\mathcal{W}\subseteq S_n$ satisfies \Cref{cond:B2} and $|\mathcal{W}|=2$, then it satisfies \Cref{cond:B3} as well.
\end{lemma}

\begin{proof}
    The previous two lemmas prove \Cref{cond:B3}(a). \Cref{cond:B3}(b) is trivial for pairs.
\end{proof}

%----------------------------------------------
\subsubsection{The proof for the general case}
%----------------------------------------------
With the case of pairs established, we move onto the general proof.


\begin{proof}[Proof that \Cref{cond:B2} implies \Cref{cond:B3}.]
    Let $\mathcal{W}= \{w_1,\dots,w_m\} \subseteq S_n$ satisfy \Cref{cond:B2}. 
    Reorder $\mathcal{W}$ lexicographically by the tuples $[i(w),j(w),-p(w),-q(w)]$ as defined in \Cref{def:Deltaw} and then relabel so that $w_i<w_{i+1}$ in that order.
    We work by induction on $m$.  

    The base case $m=1$ is trivial, and the base case $m=2$ is \Cref{lem:2 implies 3 for pairs}.
    
    Now assume that \Cref{cond:B2} implies \Cref{cond:B3} for all subsets of size at most $m-1$ where $m\geq 3$. 

    In order to utilize the induction hypothesis, we first want to show that $\{w_1,\dots,w_{m-1}\}$ satisfies \Cref{cond:B2}.
    A symmetric argument replacing each index $i$ with $m+1-i$ shows that $\{w_2,\dots,w_m\}$ also satisfies \Cref{cond:B2}.
    Putting those together with a simple induction argument, we see that $\{w_a,w_{a+1},\dots,w_{b-1},w_b\}$ satisfies \Cref{cond:B2} (and hence \Cref{cond:B3} by induction) for all $1\leq a<b\leq m$.

    Now to show that $\{w_1,\dots,w_{m-1}\}$ satisfies \Cref{cond:B2}, we must show that 
    \[
    \left(\bigcap_{i=1}^{m-1} I_{\Delta_{w_i}}\right)+I_n=\bigcap_{i=1}^{m-1} \left(I_{\Delta_{w_i}}+I_n\right)
    \]
    in the polynomial ring.
    As the kernel of the projection from the polynomial ring to $R_n$ is $I_n$ which is contained in both sides, showing that equality is equivalent to showing  
    \[
    \overline{\bigcap_{i=1}^{m-1} I_{\Delta_{w_i}}}=\bigcap_{i=1}^{m-1} \overline{I_{\Delta_{w_i}}}=\bigcap_{i=1}^{m-1}J_{w_i}.
    \]
    By hypothesis, 
    \[
    \overline{\bigcap_{i=1}^{m-1} I_{\Delta_{w_i}}}\supseteq \overline{\bigcap_{i=1}^{m} I_{\Delta_{w_i}}}=\bigcap_{i=1}^{m} J_{w_i}\subseteq \bigcap_{i=1}^{m-1} J_{w_i},
    \]
    so showing the desired equality is equivalent to showing
    \[\frac{\left(\displaystyle
\overline{\bigcap_{i=1}^{m-1}I_{\Delta_{w_i}}}\right)}
{\left(\displaystyle
\overline{\bigcap_{i=1}^{m}I_{\Delta_{w_i}}}\right)}
=
\frac{\displaystyle\bigcap_{i=1}^{m-1}J_{w_i}}
{\displaystyle\bigcap_{i=1}^{m}J_{w_i}}.\]
    As the identity is an inclusion $\bigcap_{i=1}^{m-1} I_{\Delta_{w_i}}\hookrightarrow I_{\Delta_{w_i}}$ for all $i$, it is also an injection $\overline{\bigcap_{i=1}^{m-1} I_{\Delta_{w_i}}}\hookrightarrow \bigcap_{i=1}^{m-1} \overline{I_{\Delta_{w_i}}}$, which in turn induces an injection 
    \[\frac{\left(~\displaystyle
\overline{\bigcap_{i=1}^{m-1}I_{\Delta_{w_i}}}~\right)}
{\left(~\displaystyle
\overline{\bigcap_{i=1}^{m}I_{\Delta_{w_i}}}~\right)}
\lhook\joinrel\xrightarrow{\hspace{.2in}}
\frac{\displaystyle\bigcap_{i=1}^{m-1}J_{w_i}}
{\displaystyle\bigcap_{i=1}^{m}J_{w_i}}.\]
    Thus, showing the equality is equivalent to showing that this map is surjective.
    Let $\overline{f}$ be a class in the codomain.
    As $I_{\Delta_{w_{m-1}}}$ surjects onto the numerator of the codomain, we may assume without loss of generality that $f \in I_{\Delta_{w_{m-1}}}$.
    As $I_{\Delta_{w_{m-1}}}$ is generated by squarefree monomials, if $x_H^{\vec{a}} = x_{h_1}^{a_1}\dots x_{h_K}^{a_k}$ is a monomial which appears in $f$ with nonzero coefficient, then $H\not\in \Delta_{w_{m-1}}$. 
    Let $g$ be the polynomial gotten from $f$ by setting all of the coefficients of monomials $x_H^{\vec{a}}$ to $0$ if $H \not\in \Delta_{w_m}$.
    For every $x_H^{\vec{a}}$ appearing in $g$ with nonzero coefficient, we then have $H \not\in \Delta_{w_{m-1}}$ and $H\in\Delta_{w_m}$.
    As $\Delta_{w_k}\cap \Delta_{w_m}\subseteq \Delta_{w_{m-1}}\cap \Delta_{w_m}$ by \Cref{lem:antichain ordering},  $H\not\in \Delta_{w_k}$ for all $k\in [1,m-1]$.
    Thus, $x_H\in \bigcap_{i=1}^{m-1}I_{\Delta_{w_i}}$ and therefore so is $g$.
    It follows that $\overline{g}\in \bigcap_{i=1}^{m-1}J_{w_i}$.
    As that is an ideal and (the image of) $\overline{f}$ is also in the ideal, $\overline{f-g} = \overline{f}-\overline{g}\in \bigcap_{i=1}^{m-1}J_{w_i}$.
    If $x_H^{\vec{a}}$ has nonzero coefficient in $\overline{f-g}$, by construction $H\not\in \Delta_{w_m}$.
    For such $H$, we have $x_H \in I_{\Delta_{w_m}}$.
    Thus $f-g\in I_{\Delta_{w_m}}$, so $\overline{f-g}\in\bigcap_{i=1}^m J_{w_i}$ so $\overline{f}$ and $\overline{g}$ represent the same class in the codomain, and by injectivity, in the domain.
    As $g \in \bigcap_{i=1}^{m-1}I_{\Delta_{w_i}}$ and maps to $\overline{f}$, the map is surjective. 
    Thus, we have the needed equality to show that $\{w_1,\dots,w_{m-1}\}$ satisfies \Cref{cond:B2} as desired.    

    Assume that $\dim(\Delta_{w_1})\geq \dim(\Delta_{w_m})$; if not, there is an entirely symmetric argument to what follows after replacing the indices $i$ with $m+1-i$. 
    As $\{w_1,\dots,w_{m-1}\}$ satisfies \Cref{cond:B2}, it satisfies \Cref{cond:B3} by induction. 
    Thus, it suffices to prove the parts of \Cref{cond:B3} involving $w_m$.
    (3)(b) of \Cref{cond:B3} is automatic for $w_m$ by \Cref{lem:antichain ordering}.
    Thus, it suffices to show (3)(a), i.e.,  that the intersection of $\Delta_{w_m}$ with the union of $\Delta_{w_1}$ through $\Delta_{w_{m-1}}$ is pure of dimension $\dim(\Delta_{w_m})$. 

    Let $\Gamma = \bigcup_{i=1}^{m-1}\Delta_{w_{i}}$. 
    By the argument above, $\{w_{m-1},w_m\}$ satisfies \Cref{cond:B2} and, hence, \Cref{cond:B3}.
    In particular, $\Delta_{w_m}\cap \Delta_{w_{m-1}}$ is pure (of dimension $\min(\dim(\Delta_{w_m}),\dim(\Delta_{w_{m-1}}))$).
    As $\Delta_{w_m}\cap \Gamma =\Delta_{w_m}\cap \Delta_{w_{m-1}}$, the same statement holds for it.
    $\Delta_{w_m}$ and $\Gamma$ satisfy the hypothesis of \Cref{lem:dimension-complete} by the following string of equalities which follow by standard facts about squarefree monomial ideals and the fact that $\{w_a,w_{a+1},\dots,w_{b-1},w_b\}$ satisfies \Cref{cond:B2} for all $1\leq a<b\leq m$, 
    \begin{align*}
    (I_{\Gamma}+I_n) \cap (I_{\Delta_{w_m}}+I_n) 
    &= \left(I_{\bigcup_{i=1}^{m-1}\Delta_{w_i}}+I_n\right) \cap (I_{\Delta_{w_m}}+I_n)\\
    &= \left(\left(\bigcap_{i=1}^{m-1}I_{\Delta_{w_i}}\right)+I_n\right) \cap (I_{\Delta_{w_m}}+I_n)\\
    &= \left(\bigcap_{i=1}^{m-1}\left(I_{\Delta_{w_i}}+I_n\right)\right) \cap (I_{\Delta_{w_m}}+I_n)\\
    &= \bigcap_{i=1}^{m}\left(I_{\Delta_{w_i}}+I_n\right)\\
    &= \left(\bigcap_{i=1}^{m}I_{\Delta_{w_i}}\right)+I_n\\
    &= \left(\left(\bigcap_{i=1}^{m-1}I_{\Delta_{w_i}}\right)\cap I_{\Delta_{w_m}}\right)+I_n\\
    &= \left(I_{\bigcup_{i=1}^{m-1}\Delta_{w_i}}\cap I_{\Delta_{w_m}}\right)+I_n\\
    &=(I_{\Gamma}\cap I_{\Delta_{w_m}})+I_n,
    \end{align*}
    and so  $\dim\left(\Delta_{w_m}\cap \Gamma\right) = \min(\dim(\Delta_{w_m}),  \dim(\Gamma))$ by \Cref{lem:dimension-complete}.

    As $\dim(\Gamma) \geq \dim(\Delta_{w_1})\geq \dim(\Delta_{w_m})$ by assumption, we have \[\dim\left(\Delta_{w_m}\cap \Gamma\right) = \min(\dim(\Delta_{w_m}),  \dim(\Gamma)) =\dim(\Delta_{w_m}).\]
    Thus, \Cref{cond:B3} holds for $\mathcal{W}$ as desired, and so the result follows by induction.
\end{proof}



%--------------------------------------------
\subsection{Prunable implies ideal equality}
\label{subsec:Prune_Equality}
%--------------------------------------------

We move onto the final implication of \Cref{thm:main-B}. 
We need four algebraic lemmas before our proof.
The first lemma relates the generators of $\Delta_w$ and the ideal of its intersection with $\Delta_v$.

\begin{lemma}\label{lem:mumbo-generator-jumbo} % tim - lolz
    Let $\Delta_v,\Delta_w$ be incomparable doubly-pruned complexes on $[n]$. Let $\Gamma = \Delta_v\cap \Delta_w$. If $\Gamma$ is a doubly-pruned complex that is pure of dimension $\dim(\Delta_w)$, then $I_{\Gamma} = (g_1,\dots, g_r)$ and $I_{\Delta_w} = (g_1,\dots, g_{r-1},x_B g_r)$ for some squarefree interval monomials $g_i$ such that $x_B$ is an interval squarefree monomial adjoining $g_r$ on one side.
\end{lemma}

Note, the proof of this lemma uses a few results from the next section, but they are completely independent of this entire section.

\begin{proof}   
    Suppose $\Gamma = \Delta_v\cap \Delta_w$ is doubly pruned and pure of dimension $\dim(\Delta_w)$. 
    By \Cref{lem: doubly pruned complexes are Schubert complexes}, \Cref{cor: purity of doubly pruned complexes}, and \Cref{rem: doubly pruned purity}, either $\Gamma = \Delta_u$ for some $u\in S_n(\calP)$ or $\Gamma = \mathscr{P}([i,j])$ for some $1\leq i \leq j \leq n$.
    We deal with each of these cases separately.

    First, assume that $\Gamma = \mathscr{P}([i,j])$ for some $1\leq i \leq j \leq n$, so in particular, it has dimension $j-i.$
    The complexes $\Delta_v$ and $\Delta_w$ have vertex nonfaces outside of intervals $[i_v,j_v]$ and $[i_w,j_w]$, and any vertex nonface of $\Gamma$ is a vertex nonface of $v$ or $w$.
    Thus, $[i,j] =[\max(i_v,i_w),\min(j_v,j_w)]$.
    If $[i,j] = [i_v,j_v]$, then as $\Delta_v$ has some nonvertex nonface, $\Gamma$ would also have some nonvertex nonface (or possibly smaller).
    Thus, $[i,j] \neq [i_v,j_v]$, and by symmetry, $[i,j] \neq [i_w,j_w]$.
    Up to symmetry, we may then assume that $i = i_w >i_v$ and $j = j_v <j_w$.
    Since $\Delta_w$ is pure and the same dimension as $\Gamma$, then either $[i_w,j_w] = [i,j+1]$ and its only nonvertex minimal nonface is $[q,j+1]$ or $[i_w,j_w] = [i,j+2]$ and its nonvertex minimal nonfaces are $[i,p]$ and $[q,j+2]$ with $p<q$.
    If $[i_w,j_w] = [i,j+2]$, then as $[i,p]$ is not a minimal nonface of $\Gamma$, $j<p<q<j+2$, which is a contradiction.
    If $[i_w,j_w] = [i,j+1]$, then the only nonface of $\Delta_w$ and $\Gamma$ that are different are the nonface $j+1 \in \Gamma$ and the nonface $[q,j+1]$ so the result follows in this case.
    
    Thus, we may assume that $\Gamma = \Delta_u$ for some $u\in S_n(\calP)$.
    By assumption, $\Delta_u$ is pure, and by \Cref{lem:pure-intersection}, $\Delta_w$ is pure. 
    Hence, for both $u$ and $w$, we either have $p<q$, $p =j$, or $q=i$.
    As $\Delta_u \subset \Delta_w$, $[i(u),j(u)]\subseteq [i(w),j(w)]$. 
    Similarly, as $\Delta_u \subset \Delta_v$, $[i(u),j(u)]\subseteq [i(v),j(v)]$.
    The rest of the proof is split into the nine cases define by the three possibilities for $u$ and $w$

    \medskip

    \noindent\textbf{Case 1:} $p(u)<q(u)$ and $p(w)<q(w)$.

    \medskip
    
    In this case,  $j(u)-i(u)-2 = \dim(\Delta_u) = \dim(\Delta_w) = j(w)-i(w)-2$, so $[i(u),j(u)] = [i(w),j(w)]$.
    We denote this interval as $J$ for the rest of this case.
    If $J \subseteq [i(v)+1,j(v)-1]$, then $\Delta_w\subset \Delta_v$, which is impossible by our assumption that they are incomparable. 
    Thus, either $J \subseteq [i(v),j(v)-1]$, $J \subseteq [i(v)+1,j(v)]$, or $J=[i(v),j(v)]$.
    
    If $J \subseteq [i(v),j(v)-1]$, then the vertex $j(v)$ is a minimal nonface of $\Delta_u$ and $\Delta_w$ but is a face of $\Delta_v$ so $j(u)=j(w)<j(v)$.
    If $i(v)\not\in J$, then $\Delta_w = \Delta_w$, which violates the incomparable hypothesis, so $i(v) = i(u) = i(w)$.
    Thus, $[q(v),j(v)]$ is a nonface of $\Delta_u$ that is not minimal.
    As a result, $q(u)=q(w)\geq i(w),i(v)$.
    As $i(u)=i(v)=i(w)$, $p(u) =\min(p(v),p(w))$.
    If $p(w)\leq p(v)$, then $\Delta_w\subset \Delta_v$, which is impossible again by assumption, so $p(w)>p(v)$.
    Thus, $p(u) = p(v)<p(w)\leq j(u)$.
    The minimal non-vertex nonfaces of $u$ are $[i(u),p(u)] = [i(v),p(v)]$ and $[q(u),j(u)]= [q(w),j(w)]$.
    Thus, a set of minimal generators of $I_{\Delta_w}$ is
    \[x_1,\dots,x_{i(u)-1},\qquad x_{j(u)+1},\dots,x_n,\qquad x_{[q(u),j(u)]},\qquad x_{[i(w),p(w)]}.\]
    As $x_{[i(w),p(w)]} = x_{[i(w),p(v)]}\cdot x_{[p(v)+1,p(w)]}$, the result follows in this case.
    
    The case where $[i(w),j(w)]= [i(u),j(u)] \subseteq [i(v)+1,j(v)]$ is similar.

    We are left with the case  $[i(w),j(w)]= [i(u),j(u)] = [i(v),j(v)]$, it follows that $p(u)=\min(p(v),p(w))< \max(p(v),p(w))$ and $\min(q(v),q(w))<q(u) = \max(q(v),q(w))$.
    If \newline $\min(p(v),p(w)) = p(v)$ and $\max(q(v),q(w)) = q(v)$, then $\Delta_u=\Delta_v$ which impossible.
    The symmetric statement holds for $w$.
    Thus, ($p(u) = p(v)<p(w)$ and $q(v)<q(w)=q(u)$) or ($p(u) = p(w)<p(v)$ and $q(w)<q(v)=q(u)$)
    If  $p(u) = p(v)<p(w)$ and $q(v)<q(w)=q(u)$, then the generators of $I_{\Delta_w}$ are
    \[x_1,\dots,x_{i(u)-1},\qquad x_{j(u)+1},\dots,x_n,\qquad x_{[q(u),j(u)]},\qquad x_{[i(w),p(w)]}.\]
    Again as $x_{[i(w),p(w)]} = x_{[i(w),p(v)]}\cdot x_{[p(v)+1,p(w)]}$, the result follows in this case.
    The case where $p(u) = p(w)<p(v)$ and $q(w)<q(v)=q(u)$ follows similarly.

    Thus, the result holds in the case $p(u)<q(u)$ and $p(w)<q(w)$.

    \medskip

    \noindent\textbf{Cases 2 \& 3: } $p(u)<q(u)$ and $p(w)=j(w)$ or ($p(u)<q(u)$ and $i(w)=q(w)$.

    \medskip
    
    Then  $j(u)-i(u)-2 = \dim(\Delta_u) = \dim(\Delta_w) = j(w)-i(w)-1$, so $j(u)-i(u)>j(w)-i(w)$ which is impossible as $[i(u),j(u)]\subseteq [i(w),j(w)]$.
    Thus, the result holds in these cases.

    \medskip
    
    \noindent\textbf{Case 4:} $p(u)=j(u)$ and $p(w)<q(w)$.

    \medskip
    
    In this case, $j(u)-i(u)-1 = \dim(\Delta_u) = \dim(\Delta_w) = j(w)-i(w)-2$.
    As \newline $[i(u),j(u)]\subseteq [i(w),j(w)]$, we get $[i(u),j(u)]= [i(w)+1,j(w)]$ or $[i(u),j(u)]= [i(w),j(w)-1]$.
    
    
    If $[i(u),j(u)]= [i(w)+1,j(w)]$, then the minimal generators of $I_{\Delta_w}$ are
    \[x_1,\dots,x_{i(u)-2},\qquad x_{j(u)+1},\dots,x_n,\qquad x_{[q(w),j(w)]},\qquad x_{[i(u)-1,p(w)]},\]
    while the minimal generators of $I_{\Delta_u}$ are 
    \[x_1,\dots,x_{i(u)-2},\qquad x_{j(u)+1},\dots,x_n,\qquad x_{[q(u),j(u)]},\qquad x_{i(u)-1}.\]
    It suffices to show that $x_{[q(w),j(w)]} = x_{[q(u),j(u)]}$.
    If those are not equal, then $[q(u),j(u)]$ and $i(u)-1$ must both be nonfaces of $\Delta_v$.
    Thus, $i(u)\leq i(v)$, and as $[i(u),j(u)]\subseteq [i(v),j(v)]$, $i(u)= i(v)$.
    Further, either $[q(u),j(u)] = [i(u),j(u)] = [i(v),p(v)]$ or $[q(u),j(u)] = [q(v),j(v)]$.
    If $[q(u),j(u)] = [q(v),j(v)]$, then all nonfaces of $\Delta_u$ are nonfaces of $\Delta_v$, i.e., $\Delta_v\subseteq \Delta_w$, which violates the antichain hypothesis.
    Thus, $[q(u),j(u)] = [i(u),j(u)] = [i(v),p(v)]$; if $j(u)=j(v)$, we again have a contradiction with the antichain hypothesis, so $j(v)>j(u)$.
    Thus, $i(w)+1=i(u)=q(u)=i(v)$ and $j(w) = j(u)<j(v)$.
    As $\Delta_u\subset \Delta_w$, $q(u)\geq q(w)$.
    If $q(u)>q(w)$, then $q(w)=i(w)< p(w)$, which is a contradiction.
    Thus, $q(u)=q(w)$ so $x_{[q(w),j(w)]}=x_{[q(u),j(u)]}$ as desired.

    If $[i(u),j(u)]= [i(w),j(w)-1]$, then as $p(w)<q(w)\leq j(w)-1 = j(u)=p(u)$ and $i(u)=i(w)$, $[i(w),p(w)]\subsetneq [i(u),j(u)]$, which is impossible unless $[q(u),j(u)]\subseteq [i(w),p(w)]$.
    Thus, $j(w)-1 = j(u)\leq p(w) <q(w) <j(w)$, which is impossible for integers.

    This concludes the case of $p(u)=j(u)$ and $p(w)<q(w)$.
    
    \medskip
    
    \noindent\textbf{Case 5:} $p(u)=j(u)$ and $p(w)=j(w)$. 

    \medskip
    
    In this case, $j(u)-i(u)-1 = \dim(\Delta_u) = \dim(\Delta_w) = j(w)-i(w)-1$.
    As $[i(u),j(u)] \subseteq [i(w),j(w)]$, we have $[i(u),j(u)] = [i(w),j(w)]$.
    Thus, the vertex nonfaces are the same for $\Delta_u$ and $\Delta_w$.
    They both only have one other minimal nonface, $x_{[q(u),j(u)]}$ and $x_{[q(w),j(w)]}$, respectively.
    If $q(w)\geq q(u)$, then $\Delta_w\subseteq \Delta_u$, which is impossible by the antichain assumption and the definition of $u$.
    Thus, as $q(w)<q(u)$, $x_{[q(w),j(w)]} = x_{[q(w),q(u)-1]}x_{[q(u),j(u)]}$ and so the result holds in this case.
    This concludes the case of $p(u)=j(u)$ and $p(w)=j(w)$.

    \medskip
    
    \noindent\textbf{Case 6:} $p(u)=j(u)$ and $i(w)=q(w)$. 
    
    \medskip
    
    By the same logic as the previous case, $[i(u),j(u)] = [i(w),j(w)]$.
    If $p(w)=j(w)$, then the previous case applies, so $p(w)<j(w)$.
    Then the vertex nonfaces of $\Delta_u$ and $\Delta_w$ agree, and their non-vertex minimal nonfaces are $[q(u),j(u)]$ and $[i(w),p(w)]$, respectively.
    But any nonface of $\Delta_w$ must be a nonface of $\Delta_u$ so $[i(u),p(w)] \subsetneq [i(u),j(u)]$ must be a nonface of $\Delta_u$ so we have a contradiction, so this case is impossible.

\medskip
 
 \noindent\textbf{Cases 7-9:}

\medskip

    The remaining cases of $i(u)=q(u)$ are symmetric to the previous three cases, so the result follows.
\end{proof}

Using the previous lemma, we can prove the following ideal containment.
\begin{lemma}\label{lem:colon-ideal-oscopy}
    With the same hypotheses and notation as \Cref{lem:mumbo-generator-jumbo}, letting \newline $G = (g_1,\dots, g_{r-1})$ be the corresponding ideal in $R_n$, we have
    \begin{equation}\label{eq:colon-dagger}
    (G:g_r)\subseteq G+(x_B).
    \end{equation}
\end{lemma}

\begin{proof}
 Let $\{g_1,\dots,g_r\}$ be the squarefree monomial generators of $I_{\Delta_v\cap\Delta_w}$ given by the minimal nonfaces as in \Cref{lem:mumbo-generator-jumbo} so that the generators of $I_{\Delta_w}$ are given by $\{g_1,\dots, g_{r-1},x_B g_r\}$. By the proof of \Cref{lem:mumbo-generator-jumbo}, we may order the $g_i$ so that $g_1,\dots, g_{r-2}$ are singletons.
 If $g_{r-1}$ is also a singleton, then modding out by $G$ reduces the desired statement to $\mathrm{Ann}_{R_{n-r+1}}(g_r) = (0:g_r) \subseteq (0)+x_B$, which is immediate by \Cref{lem:ann-squarefree-proofonly}.
 If $g_{r-1}$ is not a singleton, then by \Cref{lem:double pruned intersections} either:
\begin{itemize}
    \item $p(v)<p(w)$, $g_{r-1} = x_{[q(w),j(w)]}$, $g_r = x_{[i(w),p(v)]}$ and $B = [p(v)+1,p(w)]$, or
    \item $q(w)<q(v)$, $g_{r-1} = x_{[i(w),p(w)]}$, $g_r = x_{[q(v),j(w)]}$ and $B = [q(w),q(v)-1]$.
\end{itemize}
We shall prove the lemma in the first case. The proof of the second case is analogous and left to the reader. 

Quotient $R_n$ by the singleton generators $g_1,\dots, g_{r-2}$ to get a ring $R_m$, where $m= j(w)-i(w)+1$. For notational convenience, let us re-index the variables of $R_m$ to becomes $x_1,\dots, x_m$. After re-indexing, $g_r=x_{[1,p]}$, $B = [p+1,p']$  and $g_{r-1}=x_{[q, m]}$ (where $p'<q$). By 
\Cref{lem:local-endpoint-colon-proofonly}, we have
\[(g_{r-1}:g_r)\subseteq(g_{r-1},x_{[p+1,q-1]})\subseteq(g_{r-1},x_B), \]
where the second containment follows from $p'<q$. 
Lifting this containment back to $R_n$ proves \eqref{eq:colon-dagger}.
\end{proof}

We can use the previous lemma to prove the following lemma on the surjectivity of a map between Tor groups which will eventually let us show the ideal equality needed for \Cref{cond:B2}.

\begin{lemma}
\label{lem:one-sided-koszul-transfer}
Let $\Delta_v$ and $\Delta_w$ be (doubly-pruned) simplicial complexes coming from an antichain $\{v,w\}$ with $\Delta_w$ pure such that $\Delta_v\cap \Delta_w$ is a pure doubly-pruned complex of dimension $\dim(\Delta_w)$. 
Then the map
\[ \Tor_1^S(R_n,S/I_{\Delta_w})\longrightarrow\Tor_1^S(R_n,S/I_{\Delta_v\cap \Delta_w}) \]
is surjective.
\end{lemma}

\begin{proof}
The map in the statement of the lemma arises from tensoring the following short exact (Mayer-Vietoris) sequence by $R_n$:
    \[
    0\to S/I_{\Delta\cup\Delta_w}\to S/I_\Delta \oplus S/I_{\Delta_w}\to S/I_{\Delta\cap\Delta_w}\to 0.
    \]
    In particular, we get a long exact sequence on $\Tor$, part of which is
    \[
    \Tor_1^S(R_n,S/I_\Delta)\oplus \Tor_1^S(R_n,S/I_{\Delta_w})\to\Tor_1^S(R_n,S/I_{\Delta\cap\Delta_w})\to R_n/I_{\Delta\cup\Delta_w}\to R_n/I_\Delta \oplus R_n/I_{\Delta_w}.
    \]

As both $\Delta_w$ and $\Delta_w\cap \Delta_v$ are pure, their generators form a regular sequence, so the Tor groups can be computed using the Koszul complex of their respective regular sequences.
We will show surjectivity directly from this description.

Let $f_i$ be the standard coordinate basis of $(R_n)^r$. Using the same notation as in \Cref{lem:colon-ideal-oscopy}, let
\[z=\sum_{i\leq r}h_i f_i\]
be a degree-one cycle in $K_{R_n}(g_1,\ldots,g_{r-1},g_r)$, so that $\sum_{i\leq r}h_i g_i=0$ in $R_n$.
Modulo $G$, the cycle condition gives $\overline h_r \overline g_r=0$.
In the quotient, the containment in \Cref{lem:colon-ideal-oscopy} reduces to $(0:\overline{g_r}) \subseteq (\overline{x_B})$ so there exists $\overline q\in R_n/G$ such that $\overline{x_B} \,\overline{q}=\overline h_r$.
Lifting back to $R_n$, this gives
\[h_r-x_Bq=\sum_{i<r}b_i g_i \] for some $b_i$.

Define $y=\sum_{i<r}(h_i+b_i g_r)f_i+qf_r$.
Then $y$ is a cycle in $K_{R_n}(g_1,\dots, g_{r-1},x_Bg_r)$, because
\[ \sum_{i<r}(h_i+b_i g_r)g_i+qx_Bg_r = \sum_{i<r}h_ig_i+\left(\sum_{i<r}b_i g_i+qx_B\right)g_r=\sum_{i<r}h_i g_i+h_rg_r=0. \]
The image of $y$ in $K_{R_n}(g_1,\ldots,g_{r-1},g_r)$ is $\sum_{1\leq i< r} (h_i+b_ig_r)f_i + qx_Bf_r$, which differs from $z$ by 
\begin{align*}
\sum_{i<r} b_ig_rf_i + (qx_B-h_r)f_r &= \sum_{i<r} b_ig_rf_i - \sum_{i<r} b_ig_i
f_r\\
&= d\left(\sum_{i<r} b_i f_r\wedge f_i\right)
\end{align*}
a Koszul boundary.
Thus, $y$ maps to the equivalence class of $z$, so the map
\[ \Tor_1^S(R_n,S/I_{\Delta_w})\longrightarrow\Tor_1^S(R_n,S/I_{\Delta_v\cap \Delta_w}) \]
is surjective.
\end{proof}


\begin{lemma}\label{lem:deal-int-DDw}
    Let $\Delta,\Delta_w$ be simplicial complexes on $[n]$ such that $\Delta_w$ is doubly-pruned. If $\Delta\cap\Delta_w = \Delta_v\cap \Delta_w$ for some $v$ and if this intersection is pure of dimension $\dim(\Delta_w)$, then $(I_n+I_\Delta)\cap (I_n +  I_{\Delta_w}) = I_n + I_{\Delta \cup\Delta_w}$.
\end{lemma}

\begin{proof}
    We again start with the Mayer-Vietoris sequence and tensor by $R_n$ to get a long exact sequence in tor.
    Since $\Delta\cap\Delta_w = \Delta_v\cap \Delta_w$ satisfies the hypotheses of \Cref{lem:one-sided-koszul-transfer}, the map $\Tor_1^S(R_n,S/I_{\Delta_w})\to \Tor_1^S(R_n,S/I_{\Delta\cap\Delta_w})$ is surjective, hence the first map in the long exact sequence is surjective. Thus, $R_n/I_{\Delta\cup\Delta_w}\to R_n/I_\Delta \oplus R_n/I_{\Delta_w}$ is injective. Since the kernel of this map is $(I_n + I_\Delta)\cap (I_n+I_{\Delta_w})/(I_n + I_{\Delta\cup\Delta_w})$, then we have our desired equality of ideals.
\end{proof}



\begin{proof}[Proof of (3) implies (2)]
    Suppose $\mathcal{W} \subseteq S_n(\calP)$ and the elements of $\mathcal W$ admit an ordering $w_1,w_2,\dots, w_k$ satisfying the conditions in (3). We proceed by induction on $|\mathcal{W}|$. If $k=1$, then (2) is trivial. In the general case where $k>1$, then $w_1,\dots, w_{k-1}$ is an ordering that satisfies (3), so by induction
    \[
    \bigcap_{i=1}^{k-1} (I_n + I_{\Delta_{w_i}}) = I_n + \bigcap_{i=1}^{k-1}I_{\Delta_{w_i}}.
    \]
    Letting $\Delta = \bigcup_{i=1}^{k-1} \Delta_{w_i}$, then $\Delta\cap\Delta_{w_k} = \Delta_{w_t}\cap \Delta_{w_k}$ for some $t<k$ and is pure of dimension $\dim(\Delta_{w_k})$. Thus, $\Delta, \Delta_{w_k}$ satisfy the hypotheses of \Cref{lem:deal-int-DDw}, hence 
    \begin{align*}
    \bigcap_{w\in \mathcal{W}}(I_n + I_{\Delta_{w}}) &= \bigcap_{i=1}^{k-1} (I_n + I_{\Delta_{w_i}})\cap(I_n + I_{\Delta_{w_k}})\\
    &= (I_n + \bigcap_{i=1}^{k-1}I_{\Delta_{w_i}})\cap(I_n + I_{\Delta_{w_k}})\\
    &= (I_n + I_{\Delta})\cap(I_n + I_{\Delta_{w_k}})\\
    &= I_n + I_{\Delta\cup\Delta_{w_k}} = I_n + \bigcap_{w\in \mathcal{W}}I_{\Delta_w}\qedhere
    \end{align*}
\end{proof}

%----------------------------------------------------------------------
\subsection{The proof of \texorpdfstring{\Cref{thm:main-A}}{Theorem A}}
%----------------------------------------------------------------------
As we have proved \Cref{thm:main-B}, we can now complete the proof of \Cref{thm:main-A}.
\begin{proof}[Proof of \Cref{thm:main-A}]

    The equivalence of \Cref{cond:A1,cond:A2} is exactly the equivalence of \Cref{cond:B1,cond:B2} when $\mathcal{W} = \{w\}$ as then $J_\mathcal{W} = J_w$ and the equality in \Cref{cond:A2} is trivial.
    Then \Cref{lem: A3 implies A1} and \Cref{lem: A2 implies A3} give the remaining needed implications of \Cref{cond:A3} implies  \Cref{cond:A1} and \Cref{cond:A2} implies \Cref{cond:A3}, respectively. 
\end{proof}